% ==============================================================================
% 07_anhang.tex – Anhang (wird nach dem Literaturverzeichnis eingebunden)
% ==============================================================================
\appendix
\chapter{Anhang}
\label{ch:anhang}

\section{Befehle zum Wiederholen der Experimente}
\label{sec:anhang_befehle}
Alle Befehle werden im Hauptordner des Projekts ausgeführt (siehe auch
\code{README.md}).

\paragraph{Daten vorbereiten.}
\begin{verbatim}
python main.py
\end{verbatim}

\paragraph{Baselines und Hyperparameterstudie (GazeCapture).}
\begin{verbatim}
python exp.py --experiment baseline
python exp.py --experiment all --sample-seed 0 --val-size 200
python exp.py --experiment trainable_layers --include-resnet
\end{verbatim}

\paragraph{Räumliche Analyse und Abbildungen.}
\begin{verbatim}
python spatial_analysis.py --input results/heatmap/base_loss1.csv \
       --compare results/heatmap/base_loss2.csv --out results/figures
python plot_targets.py --out results/figures
\end{verbatim}

\paragraph{Personalisierung (E0--E3).}
\begin{verbatim}
python train_base.py
python personal_exp.py --sessions all --base norm_subject norm_random
python personal_exp.py --sessions gen4 --labels clean
python personal_exp.py --sessions 21 22 --split random --seeds 42 43 44
python personal_exp.py --sessions 21 22 --split sequential --seeds 42 43 44
python personal_exp.py --split two_videos --pairs 21:22,22:21 --seeds 42 43 44
python personal_exp.py --sessions all --split sequential
python personal_exp.py --summary
\end{verbatim}

\paragraph{Gestaltung der Sitzungen und bereinigte Labels.}
\begin{verbatim}
python session_design.py
python session_design.py --write-clean
\end{verbatim}

Einfacher geht es mit dem Notebook \code{notebooks/Masterarbeit.ipynb}: Zelle~1 trainiert
die Basismodelle (E0), Zelle~3 führt die Personalisierung (E1 bis E3) aus und
erstellt alle Tabellen.
